\section{Related Work}
\label{sec:related}

\para{LLM steganographic encoding.}
The most directly relevant line of work evaluates how well LLMs can produce steganographic text when prompted. \citet{karpov2026nest} introduced NEST, a systematic evaluation of acrostic steganography across 28 models and 10 difficulty levels. They showed that Claude Opus 4.5 achieves 92\% per-number accuracy at sequence length $D{=}4$ but drops to 12\% at $D{=}50$, establishing that acrostic accuracy degrades with message length. \citet{zolkowski2025early} extended this by comparing two encoding schemes---\texttt{sentence\_letter} (acrostic) and \texttt{third\_word\_letter}---finding that models perform substantially worse on the latter. They also showed that models default to acrostic encoding when asked to steganographically encode without specifying a scheme. \citet{karpov2025potentials} explored reinforcement learning for steganography, showing that RL fine-tuning improves covert communication by roughly 20\% for 3 bits of entropy, and that performance declines 33.5\% on average as payload length increases. Unlike these works, we systematically vary dilution level across six schemes and measure encoding, extraction, and detection in a unified framework.

\para{Steganographic collusion and AI safety.}
A parallel line of research examines steganography as a safety concern in multi-agent systems. \citet{motwani2024secret} formalized secret collusion among AI agents, testing GPT-4 for steganographic message passing and finding up to 26\% success rate in prompted settings. \citet{mathew2024hidden} demonstrated that steganographic collusion can emerge from misspecified reward incentives during training, not just from prompting, and found that standard mitigations such as paraphrasing are insufficient. \citet{roger2023preventing} showed that language models can be fine-tuned to hide reasoning in chain-of-thought outputs and proposed paraphrasing as a defense, which limits encoding to roughly 3 bits per kilobyte. These works establish the safety motivation for understanding steganographic capabilities but do not characterize how those capabilities vary with signal dilution---the central question of our study.

\para{Classical text steganography.}
Before the LLM era, text steganography relied on algorithmic methods such as synonym substitution, syntactic transformation, and format-based encoding~\citep{ziegler2019neural}. \citet{ziegler2019neural} proposed neural linguistic steganography using arithmetic coding with language models, achieving high capacity with low perplexity overhead. These methods differ fundamentally from prompted LLM steganography: they use the model as a tool for sampling steganographic text rather than relying on the model to understand and implement an encoding scheme. Our work focuses on the latter setting, which is more relevant to AI safety concerns about models spontaneously developing or following steganographic instructions.
